% ==============================================================================
% 公共题库: teaching/pool/解三角形/难题/q_triangle_dof_range_parameterization.tex
% 难度: ★★★ 难题
% 知识点: 剩余自由度1(DOF=1)、圆弧滑动轨迹、单变量消元参数化、自变量定义域裁剪
% ==============================================================================

\newcommand{\qTriangleDofRangeStem}{%
在 \(\triangle ABC\) 中，角 \(A, B, C\) 的对边分别为 \(a, b, c\)。已知 \(C = \frac{\pi}{3}\)，\(c = \sqrt{3}\)。
\begin{enumerate}[(1)]
  \item 求 \(\triangle ABC\) 外接圆半径 \(R\)；
  \item 求 \(\triangle ABC\) 的周长 \(L = a+b+c\) 的取值范围。
\end{enumerate}%
}

\newcommand{\qTriangleDofRangeAnswer}{%
(1) \(R = 1\)；\qquad (2) 周长 \(L \in (2\sqrt{3}, 3\sqrt{3}]\)。
}

\newcommand{\qTriangleDofRangeSolution}{%
(1) \(2R = c / \sin C = 2 \implies R = 1\)。
(2) 消元参数化：设 \(A\) 为参数，\(A \in (0, \frac{2\pi}{3})\)。
\(a=2\sin A\), \(b=2\sin(\frac{2\pi}{3}-A)\)。
\(L(A) = 2\sqrt{3}\sin(A+\frac{\pi}{6})+\sqrt{3}\)。
因 \(A+\frac{\pi}{6} \in (\frac{\pi}{6}, \frac{5\pi}{6})\)，最大值为 \(3\sqrt{3}\)（在 \(A=\frac{\pi}{3}\) 取得），下界趋近于 \(2\sqrt{3}\)。
故周长范围为 \((2\sqrt{3}, 3\sqrt{3}]\)。
}

\newcommand{\qTriangleDofRangeDiagnosis}{%
【错因】未导出自变量区间 $A \in (0, \frac{2\pi}{3})$ 导致无法求出下界，或将开端点错写为闭端点。
}
